\documentclass[10pt,a4paper]{amsart}
\usepackage{amsmath,amssymb,amsthm,graphicx}
% Retained minimal preamble for the Stacks Project chapter more-algebra.tex (More on Algebra).
% The chapter file pinned-original-source.tex carries only \input{preamble} on its first line, so no
% preamble text lives inside the chapter itself.  The Stacks Project uses one single shared preamble for
% every chapter; the snapshot stages it identically as sources/stacks-chow/preamble.tex and
% sources/stacks-etale-cohomology/preamble.tex.  The block below is copied mechanically, line for line,
% from sources/stacks-chow/preamble.tex lines 14-18 (the Xy-pic setup) and 154-217 (the native theorem
% declarations together with the mathematical macros).  Nothing was added, removed or redefined; the
% wrapper class and the amsmath/amssymb/amsthm/graphicx base packages are supplied by the assembler.
% Chapter: more-algebra.tex -- More on Algebra.
% For commutative diagrams we use Xy-pic
\usepackage[all]{xy}

% We use 2cell for 2-commutative diagrams.
\xyoption{2cell}

% Theorem environments.
%
\theoremstyle{plain}
\newtheorem{theorem}[subsection]{Theorem}
\newtheorem{proposition}[subsection]{Proposition}
\newtheorem{lemma}[subsection]{Lemma}

\theoremstyle{definition}
\newtheorem{definition}[subsection]{Definition}
\newtheorem{example}[subsection]{Example}
\newtheorem{exercise}[subsection]{Exercise}
\newtheorem{situation}[subsection]{Situation}

\theoremstyle{remark}
\newtheorem{remark}[subsection]{Remark}
\newtheorem{remarks}[subsection]{Remarks}

\numberwithin{equation}{subsection}

% Macros
%
\def\lim{\mathop{\mathrm{lim}}\nolimits}
\def\colim{\mathop{\mathrm{colim}}\nolimits}
\def\Spec{\mathop{\mathrm{Spec}}}
\def\Hom{\mathop{\mathrm{Hom}}\nolimits}
\def\Ext{\mathop{\mathrm{Ext}}\nolimits}
\def\SheafHom{\mathop{\mathcal{H}\!\mathit{om}}\nolimits}
\def\SheafExt{\mathop{\mathcal{E}\!\mathit{xt}}\nolimits}
\def\Sch{\mathit{Sch}}
\def\Mor{\mathop{\mathrm{Mor}}\nolimits}
\def\Ob{\mathop{\mathrm{Ob}}\nolimits}
\def\Sh{\mathop{\mathit{Sh}}\nolimits}
\def\NL{\mathop{N\!L}\nolimits}
\def\CH{\mathop{\mathrm{CH}}\nolimits}
\def\proetale{{pro\text{-}\acute{e}tale}}
\def\etale{{\acute{e}tale}}
\def\QCoh{\mathit{QCoh}}
\def\Ker{\mathop{\mathrm{Ker}}}
\def\Im{\mathop{\mathrm{Im}}}
\def\Coker{\mathop{\mathrm{Coker}}}
\def\Coim{\mathop{\mathrm{Coim}}}

% Boxtimes
%
\DeclareMathSymbol{\boxtimes}{\mathbin}{AMSa}{"02}

%
% Macros for moduli stacks/spaces
%
\def\QCohstack{\mathcal{QC}\!\mathit{oh}}
\def\Cohstack{\mathcal{C}\!\mathit{oh}}
\def\Spacesstack{\mathcal{S}\!\mathit{paces}}
\def\Quotfunctor{\mathrm{Quot}}
\def\Hilbfunctor{\mathrm{Hilb}}
\def\Curvesstack{\mathcal{C}\!\mathit{urves}}
\def\Polarizedstack{\mathcal{P}\!\mathit{olarized}}
\def\Complexesstack{\mathcal{C}\!\mathit{omplexes}}
% \Pic is the operator that assigns to X its picard group, usage \Pic(X)
% \Picardstack_{X/B} denotes the Picard stack of X over B
% \Picardfunctor_{X/B} denotes the Picard functor of X over B
\def\Pic{\mathop{\mathrm{Pic}}\nolimits}
\def\Picardstack{\mathcal{P}\!\mathit{ic}}
\def\Picardfunctor{\mathrm{Pic}}
\def\Deformationcategory{\mathcal{D}\!\mathit{ef}}
\usepackage{cleveref}
\crefname{originallemma}{original lemma}{original lemmas}
\Crefname{originallemma}{Original lemma}{Original lemmas}
\makeatletter
\newlabel{more-algebra-lemma-p-basis}{{07P2}{source}{}{sourcepointer-1}{}}
\newlabel{more-algebra-lemma-p-basis@cref}{{[originallemma][0][]07P2}{[1][0][]source}{}{}{}}
\makeatother
\title{p-independent sets, p-bases and the differentials of a field extension of characteristic p}
\author{The Stacks Project Authors}
\date{stacks-project:more-algebra.tex:local-snapshot-20261012; GNU FDL 1.2 or later}
\begin{document}
\maketitle
\setcounter{subsection}{1}
\noindent\textit{Editorial note.} All delivered mathematical text is the author's own, copied from the pinned Stacks Project chapter file more-algebra.tex; no statement or proof step was rewritten, reordered or removed.\par\medskip
\noindent\textit{Editorial note.} The chapter file pinned-original-source.tex carries only \textbackslash{}input\{preamble\} on its first line, so no preamble text lives inside the chapter. The retained file sources/stacks-more-algebra/retained-preamble.tex is a mechanical, line-for-line copy of the project's single shared preamble as staged in the snapshot, namely sources/stacks-chow/preamble.tex lines 14-18 (Xy-pic setup) and 154-217 (theorem declarations and mathematical macros). Nothing was added, removed or redefined, no package was added, and the excerpt serves only as preamble, never as a span coordinate.\par\medskip
\noindent\textit{Editorial note.} Both counted spans start at the \textbackslash{}begin line of their environment, so the statement and the author's proof are delivered whole; surrounding prose of the section is deliberately outside the counted spans and is carried as standard prerequisites instead.\par\medskip
\noindent\textit{Editorial note.} The citation fields (source\_id, source\_version, authors, source\_url, license.name, source\_title) are identical to those of the cards already published from this chapter, so that the batch selector treats them as one source instead of deferring them as a same-source clash.\par\medskip
\noindent\textit{Allowed standard prerequisites (editorial).}\begin{itemize}
\item p-independence and p-basis of a field extension, as defined immediately before the statement. definition environment at source lines 11338-11358 of the same chapter. Supplies the two properties that are being shown equivalent and the monomials x\textasciicircum{}E with exponents below p.
\item Kaehler differentials of a field extension and the fact that dx\_i form a K-basis when K is generated by the x\_i over kK\textasciicircum{}p. used at proof lines 11394-11398 to derive (b) from (a). Provides the module in which the differentials live and the presentation K = kK\textasciicircum{}p[X\_i]/(X\_i\textasciicircum{}p - x\_i\textasciicircum{}p) that makes the computation possible.
\item Existence of k-derivations dual to a p-independent set, and the fact that every k-derivation kills kK\textasciicircum{}p. used at proof lines 11381-11387, announced inside the proof as theta\_i with theta\_i(x\_j) = delta\_ij. Produces the new relations that lead to the degree contradiction, which is the whole content of the first implication.
\end{itemize}\medskip
\ifdefined\hypertarget\hypertarget{sourcepointer-1}{}\fi
\noindent\textit{Source reference.} Original lemma 07P2: focused-seeds/AGP/sources/stacks-snapshot-20261011/tags.txt, source lines 9782--9782. The lemma's own canonical tag, quoted so that readers can resolve the label inside the delivered text; it points at the counted statement itself and carries no proof content.\par\medskip
\begin{lemma}
\label{lemma-p-basis}
Let $K/k$ be a field extension. Assume $k$ has characteristic
$p > 0$. Let $\{x_i\}$ be a subset of $K$. The following are equivalent
\begin{enumerate}
\item the elements $\{x_i\}$ are $p$-independent over $k$, and
\item the elements $\text{d}x_i$ are $K$-linearly independent
in $\Omega_{K/k}$.
\end{enumerate}
Any $p$-independent collection can be extended to a $p$-basis of $K$ over $k$.
In particular, the field $K$ has a $p$-basis over $k$.
Moreover, the following are equivalent:
\begin{enumerate}
\item[(a)] $\{x_i\}$ is a $p$-basis of $K$ over $k$, and
\item[(b)] $\text{d}x_i$ is a basis of the $K$-vector space $\Omega_{K/k}$.
\end{enumerate}
\end{lemma}

\begin{proof}
Assume (2) and suppose that $\sum a_E x^E = 0$ is a linear relation
with $a_E \in k K^p$. Let $\theta_i : K \to K$ be a $k$-derivation such that
$\theta_i(x_j) = \delta_{ij}$ (Kronecker delta). Note that any $k$-derivation
of $K$ annihilates $kK^p$. Applying $\theta_i$ to the given relation we
obtain new relations
$$
\sum\nolimits_{E, e_i > 0}
e_i a_E x_1^{e_1}\ldots x_i^{e_i - 1} \ldots x_n^{e_n} = 0
$$
Hence if we pick $\sum a_E x^E$ as the relation with minimal
total degree $|E| = \sum e_i$ for some $a_E \not = 0$, then we
get a contradiction. Hence (1) holds.

\medskip\noindent
If $\{x_i\}$ is a $p$-basis for $K$ over $k$, then
$K \cong kK^p[X_i]/(X_i^p - x_i^p)$. Hence we see that
$\text{d}x_i$ forms a basis for $\Omega_{K/k}$ over $K$.
Thus (a) implies (b).

\medskip\noindent
Let $\{x_i\}$ be a $p$-independent subset of $K$ over $k$. An application
of Zorn's lemma shows that we can enlarge this to a maximal $p$-independent
subset of $K$ over $k$. We claim that any maximal $p$-independent subset
$\{x_i\}$ of $K$ is a $p$-basis of $K$ over $k$. The claim will imply
that (1) implies (2) and establish the existence of $p$-bases.
To prove the claim let $L$ be the subfield of $K$ generated by
$kK^p$ and the $x_i$. We have to show that $L = K$. If $x \in K$
but $x \not \in L$, then $x^p \in L$ and $L(x) \cong L[z]/(z^p - x)$.
Hence $\{x_i\} \cup \{x\}$ is $p$-independent over $k$, a contradiction.

\medskip\noindent
Finally, we have to show that (b) implies (a). By the equivalence of (1)
and (2) we see that $\{x_i\}$ is a maximal $p$-independent subset
of $K$ over $k$. Hence by the claim above it is a $p$-basis.
\end{proof}

\end{document}
